\par
\SourceLeaf{172}
\textit{Ex.} Sit $D=79$ eruntque valores ipsius $a$ viginti duo ${\mp1}$, $2$, $3$, $5$, $6$, $7$, $9$, $10$, $13$, $14$, $15$. Unde inveniuntur formae undeviginti:
\[
\begin{gathered}
(1,8,-15),\quad (2,7,-15),\quad {(3,8,-5)},\quad (3,7,-10),\\
(5,8,-3),\quad (5,7,-6),\quad (6,7,-5),\quad (6,5,-9),\\
(7,4,-9),\quad (7,3,-10),\quad (9,5,-6),\quad (9,4,-7),\\
(10,7,-3),\quad (10,3,-7),\quad (13,1,-6),\quad (14,3,-5),\\
(15,8,-1),\quad (15,7,-2),\quad (15,2,-5)
\end{gathered}
\]
totidemque aliae quae fiunt ex his, si terminorum exterorum signa {commutantur. puta} $(-1,8,15)$, $(-2,7,15)$ etc., ita ut omnes triginta octo sint. Sed ex his reiiciendae sex $(\pm13,1,\mp6)$, $(\pm14,3,\mp5)$, $(\pm15,2,\mp5)$, reliquae triginta duae omnes reductas amplectuntur. Per methodum secundam eaedem formae prodeunt sequenti ordine\HistoricalNote{*)}{Pro $b=1$, ${-78}$ in duos factores, qui neglecto signo inter $\surd79+1$ et $\surd79-1$ iaceant, resolvi nequit; quare hic valor est praetereundus, ex eademque ratione valores $2$ et $6$.}:
\[
\begin{gathered}
(\pm7,3,\mp10),\quad (\pm10,3,\mp7),\quad {\pm7,4,\mp9)},\quad (\pm9,4,\mp7),\\
(\pm6,5,\mp9),\quad (\pm9,5,\mp6),\quad (\pm2,7,\mp15),\quad (\pm3,7,\mp10),\\
(\pm5,7,\mp6),\quad (\pm6,7,\mp5),\quad (\pm10,7,\mp3),\quad (\pm15,7,\mp2),\\
(\pm1,8,\mp15),\quad (\pm3,8,\mp5),\quad (\pm5,8,\mp3),\quad (\pm15,8,\mp1).
\end{gathered}
\]

\Needspace{8\baselineskip}
\Article{186}
Sit $F$ forma reducta determinantis $D$, ipsique ab ultima parte contigua forma reducta $F'$; huic iterum ab ultima parte contigua reducta $F''$; reducta $F'''$ ipsi $F''$ contigua ab ultima parte etc. Tum patet, omnes formas $F',F'',F'''$ etc. esse prorsus determinatas, et tum inter se tum formae $F$ proprie aequivalentes. Quoniam vero multitudo omnium formarum reductarum determinantis dati est finita, manifestum est, omnes formas in progressione infinita $F,F',F''$ etc. diversas esse non posse. Ponamus $F^m$ et $F^{m+n}$ esse identicas, eruntque $F^{m-1}$, $F^{m+n-1}$ reductae, eidem formae reductae a parte prima contiguae, adeoque identicae; hinc eodem modo $F^{m-2}$ et $F^{m+n-2}$ etc. tandemque $F$ et $F^n$ identicae erunt. Quare in progressione $F,F',F''$ etc., si modo satis longe continuatur, necessario tandem forma prima $F$ recurret; et si supponimus $F^n$ esse primam identicam cum $F$, sive omnes $F',F''\mathinner{\ldotp\ldotp\ldotp\ldotp} F^{n-1}$ a forma $F$ diversas: facile perspicitur, omnes formas $F,F',F''\mathinner{\ldotp\ldotp\ldotp\ldotp} F^{n-1}$ diversas fore. Complexum harum for%
\SourceLeaf{173}%
marum vocabimus \textit{periodum formae} $F$. Si igitur progressio ultra ultimam periodi formam producitur, eaedem formae $F,F',F''$ etc. iterum prodibunt, progressioque tota infinita $F,F',F''$ etc. constituta erit ex hac periodo formae $F$ infinities repetita.

Progressio $F,F',F''$ etc. etiam retro continuari potest, praeponendo formae $F$ reductam ${}'F$, quae ipsi a parte prima est contigua; huic iterum reductam ${}''F$, quae ipsi a prima parte contigua etc. Hoc modo habebitur progressio formarum \textit{utrimque} infinita
\[
\mathinner{\ldotp\ldotp\ldotp\ldotp} {}'''F,\ {}''F,\ {}'F,\ F,\ F',\ F'',\ F'''\mathinner{\ldotp\ldotp\ldotp\ldotp}
\]
perspicieturque facile, ${}'F$ identicam fore cum $F^{n-1}$, ${}''F$ cum $F^{n-2}$ etc. adeoque progressionem etiam a laeva parte e periodo formae $F$, infinities repetita, esse constitutam.

Si formis $F,F',F''$ etc. ${}'F$, ${}''F$ etc. tribuuntur indices $0$, $1$, $2$ etc. $-1$, $-2$ etc. generaliterque formae $F^m$ index $m$, formae ${}^mF$ index $-m$, patet, \textit{formas quascunque seriei identicas fore vel diversas, prout ipsarum indices congrui sint vel incongrui secundum modulum $n$.}

\textit{Ex.} Periodus formae $(3,8,-5)$, cuius determinans $=79$, invenitur haec: ${3,8,-5)}$, $(-5,7,6)$, $(6,5,-9)$, $(-9,4,7)$, $(7,3,-10)$, $(-10,7,3)$. Post ultimam iterum prodit $(3,8,-5)$. Hic itaque $n=6$.

\Needspace{8\baselineskip}
\Article{187}
Ecce quasdam observationes generales circa has periodos.

1) Si formae $F,F',F''$ etc.; ${}'F,{}''F,{}'''F$ etc. ita exhibentur:
\[
\begin{gathered}
{a,b,-a')},\quad (-a',b',a''),\quad (a'',b'',-a''')\text{ etc.};\\
(-{}'a,{}'b,a),\quad ({}''a,{}''b,-{}'a),\quad (-{}'''a,{}'''b,{}''a)
\end{gathered}
\]
omnes $a,a',a'',a'''$ etc. ${}'a,{}''a,{}'''a$ etc. \textit{eadem signa} habebunt (art.~184,~1), omnes vero $b,b',b''$ etc. ${}'b,{}''b$ etc. erunt positivi.

2) Hinc manifestum est, numerum $n$ (multitudinem formarum ex quibus periodus formae $F$ constat) semper esse \textit{parem}. Etenim terminus primus formae cuiusvis $F^m$ ex hac periodo manifesto idem signum habebit uti terminus primus $a$ formae $F$, si $m$ est par, oppositum, si $m$ est impar. Quare quum $F^n$ et $F$ identicae sint, $n$ necessario erit par.
\SourceLeaf{174}

\Needspace{14\baselineskip}
3) Algorithmus, per quem numeri $b',b'',b'''$ etc., $a'',a'''$ etc. inveniuntur, ex art.~184,~6 est hic:
\[
\begin{array}{lll}
b'\equiv-b\ (\mathrm{mod.}\ a')&\text{inter limites }\surd D\text{ et }\surd D\mp a';&a''=\dfrac{D-b'b'}{a'}\\[6pt]
b''\equiv-b'\ (\mathrm{mod.}\ a'')&\mathinner{\ldotp\ldotp\ldotp\ldotp\ldotp\ldotp\ldotp\ldotp\ldotp}\quad\surd D\mp a'';&a'''=\dfrac{D-b''b''}{a''}\\[6pt]
b'''\equiv-b''\ (\mathrm{mod.}\ a''')&\mathinner{\ldotp\ldotp\ldotp\ldotp\ldotp\ldotp\ldotp\ldotp\ldotp}\quad\surd D\mp a''';&a''''=\dfrac{D-b'''b'''}{a'''}
\end{array}
\]
\centerline{etc.}
ubi in columna secunda signa superiora vel inferiora sunt accipienda, prout $a,a',a''$ etc. sunt positivi vel negativi. Loco formularum in columna tertia etiam sequentes adhiberi possunt, quae commodiores evadunt, quando $D$ est numerus magnus:
\[
\begin{aligned}
a''&=\frac{b+b'}{a'}(b-b')+a\\
a'''&=\frac{b'+b''}{a''}(b'-b'')+a'\\
a''''&=\frac{b''+b'''}{a'''}({b''}-b''')+a''
\end{aligned}
\]
\centerline{etc.}

4) Forma quaecunque $F^m$, in periodo formae $F$ contenta, proprie eandem periodum habet ut $F$. Scilicet periodus illa erit $F^m$, $F^{m+1}\mathinner{\ldotp\ldotp\ldotp\ldotp} F^{n-1}$, $F,F'\mathinner{\ldotp\ldotp\ldotp\ldotp} F^{m-1}$, in qua eaedem formae eodemque ordine occurrunt, ut in periodo formae $F$, et quae ab hac tantummodo respectu initii et finis discrepat.

5) Hinc patet, omnes formas reductas eiusdem determinantis $D$ in periodos \textit{distribui} posse. Accipiatur aliqua harum formarum, $F$, ad libitum investigeturque ipsius periodus, $F,F',F''\mathinner{\ldotp\ldotp\ldotp\ldotp} F^{n-1}$, quam designemus per $P$. Si haec omnes formas reductas determinantis $D$ nondum amplectitur, sit aliqua in ipsa non contenta $G$ huiusque periodus $Q$. Tum patet $P$ et $Q$ nullam formam communem habere posse; alioquin enim etiam $G$ in $P$ contenta esse deberet periodique omnino coinciderent. Si $P$ et $Q$ omnes formas reductas nondum exhauriunt, aliqua ex deficientibus, $H$, periodum tertiam, $R$, suppeditabit, quae neque cum $P$ neque cum $Q$ formam communem habebit. Hoc modo continuare possumus, usquedum omnes formae reductae sint exhaustae. Ita \textit{e. g.} omnes formae reductae determinantis $79$ in sex periodos distribuuntur:
\SourceLeaf{175}
\[
\begin{array}{rl}
\mathrm{I.}&(1,8,-15),\quad (-15,7,2),\quad (2,7,-15),\quad (-15,8,1).\\
\mathrm{II.}&(-1,8,15),\quad (15,7,-2),\quad (-2,7,15),\quad (15,8,-1).\\
\mathrm{III.}&(3,8,-5),\quad (-5,7,6),\quad (6,5,-9),\\
&(-9,4,7),\quad (7,3,-10),\quad (-10,7,3).\\
\mathrm{IV.}&(-3,8,5),\quad (5,7,-6),\quad (-6,5,9),\\
&(9,4,-7),\quad (-7,3,10),\quad (10,7,-3).\\
\mathrm{V.}&(5,8,-3),\quad (-3,7,10),\quad (10,3,-7),\\
&(-7,4,9),\quad (9,5,-6),\quad (-6,7,5).\\
\mathrm{VI.}&(-5,8,3),\quad (3,7,-10),\quad (-10,3,7),\\
&(7,4,-9),\quad (-9,5,6),\quad (6,7,-5).
\end{array}
\]

6) Vocemus \textit{formas socias}, quae ex iisdem terminis constant, sed ordine inverso positis, ut $(a,b,-a')$, $(-a',b,a)$. Tum facile perspicitur ex art.~184,~7, si periodus formae reductae $F$ sit $F,F',F''\mathinner{\ldotp\ldotp\ldotp\ldotp} F^{n-1}$, formae $F$ socia $f$ formisque $F^{n-1},F^{n-2}\mathinner{\ldotp\ldotp\ldotp\ldotp} F'',F'$ resp. sociae sint formae $f',f''\mathinner{\ldotp\ldotp\ldotp\ldotp} f^{n-2},f^{n-1}$: periodum formae $f$ fore $f,f',f''\mathinner{\ldotp\ldotp\ldotp\ldotp} f^{n-2},f^{n-1}$, adeoque ex totidem formis constare, ut periodum formae $F$. Periodos formarum sociarum vocabimus \textit{periodos socias}. Ita in exemplo nostro sociae sunt periodi III et VI; IV et V.

7) Sed fieri etiam potest, ut forma $f$ ipsa in periodo sociae suae $F$ occurrat, uti in ex. nostro in periodo I et II, adeoque periodus formae $F$ cum periodo formae $f$ conveniat, sive ut \textit{periodus formae $F$ sibi ipsi sit socia}. Quoties hoc evenit, in hac periodo duae formae ancipites invenientur. Ponamus enim periodum formae $F$ constare e $2n$ formis sive $F$ et $F^{2n}$ esse identicas; porro sit $2m+1$ index formae $f$ in periodo formae $F$\HistoricalNote{*)}{Index hic necessario erit impar, quia manifesto termini primi formarum ${F},f$ signa opposita habent (\textit{vid. supra}, $2$).}, sive $F^{2m+1}$ et $F$ sociae. Tum patet etiam $F'$ et $F^{2m}$ fore socias nec non $F''$ et $F^{2m-1}$ etc., adeoque etiam $F^m$ et $F^{m+1}$. Sit $F^m=(a^m,b^m,-a^{m+1})$, $F^{m+1}=(-a^{m+1},b^{m+1},a^{m+2})$. Tum erit $b^m+b^{m+1}\equiv0\ (\mathrm{mod.}\ a^{m+1})$; ex defin. formarum sociarum vero erit $b^m=b^{m+1}$ atque hinc $2b^{m+1}\equiv0\ (\mathrm{mod}\ a^{m+1})$, sive forma $F^{m+1}$ anceps. --- Eodem modo $F^{2m+1}$ et $F^{2n}$ erunt sociae; hinc $F^{2m+2}$ et $F^{2n-1}$; $F^{2m+3}$ et $F^{2n-2}$ etc. tandemque $F^{m+n}$ et $F^{m+n+1}$, quarum posterior erit anceps, uti per simile ratiocinium facile probatur. Quia vero $m+1$ et $m+n+1$ secundum mod.~$2n$ sunt incongrui, formae $F^{m+1}$ et $F^{m+n+1}$ identicae non erunt (art.~186, ubi $n$ idem denotat, quod hic $2n$). Ita in I sunt formae ancipites $(1,8,-15)$, $(2,7,-15)$, in II vero $(-1,8,15)$, $(-2,7,15)$.
\SourceLeaf{176}

8) Vice versa, \textit{quaevis periodus, in qua forma anceps occurrit, sibi ipsi socia est.} Facile enim perspicitur, si $F^m$ sit forma reducta anceps: formam ipsi sociam (quae etiam est reducta) simul ipsi a parte prima contiguam esse, \textit{i. e.} $F^{m-1}$ et $F^m$ socias. Tum vero tota periodus sibi ipsi socia erit. --- {Hinc patet.} \textit{fieri non posse, ut unica tantum forma anceps in periodo aliqua contenta sit.}

9) Sed etiam \textit{plures quam duae in eadem periodo esse nequeunt.} Ponamus enim in periodo formae $F$, ex $2n$ formis constante, tres formas ancipites dari $F^\lambda,F^\mu,F^\nu$, ad indices $\lambda,\mu,\nu$ respective pertinentes, ita ut $\lambda,\mu,\nu$ sint numeri inaequales inter limites $0$ et $2n-1$ (incl.) siti. Tum formae $F^{\lambda-1}$ et $F^\lambda$ erunt sociae; similiterque $F^{\lambda-2}$ et $F^{\lambda+1}$ etc. tandemque $F$ et $F^{2\lambda-1}$. Ex eadem ratione $F$ et $F^{2\mu-1}$ sociae erunt, nec non $F$ et $F^{2\nu-1}$; quare $F^{2\lambda-1}$, $F^{2\mu-1}$, $F^{2\nu-1}$ identicae, indicesque $2\lambda-1$, $2\mu-1$, $2\nu-1$ secundum modulum $2n$ congrui erunt, et proin etiam $\lambda\equiv\mu\equiv\nu\ (\mathrm{mod.}\ n)$. \textit{Q. E. A.} quia manifesto inter limites $0$ et $2n-1$ tres numeri diversi secundum modulum $n$ congrui iacere nequeunt.

\Needspace{8\baselineskip}
\Article{188}
Quum omnes formae ex eadem periodo proprie sint aequivalentes, quaestio oritur, annon etiam formae e periodis diversis proprie aequivalentes esse possint. Sed antequam ostendamus, \textit{hoc esse impossibile}, quaedam de transformatione formarum reductarum sunt exponenda.

Quoniam in sequentibus de formarum transformationibus persaepe agendum erit: ut prolixitatem quantum fieri potest evitemus, sequenti scribendi compendio abhinc semper utemur. Si forma aliqua $LXX+2MXY+NYY$ per substitutionem $X=\alpha x+\beta y$, $Y=\gamma x+\delta y$ in formam $lxx+2mxy+nyy$ transformatur: simpliciter dicemus, $(L,M,N)$ transformari in $(l,m,n)$ per substitutionem $\alpha,\beta,\gamma,\delta$. Hoc modo opus non erit, indeterminatas formarum singularum, de quibus agitur, per signa propria denotare. --- Palam vero est, indeterminatam \textit{primam} a \textit{secunda} in quavis forma probe distingui debere.

Proposita sit forma reducta $(a,b,-a')\mathinner{\ldotp\ldotp\ldotp} f$, determinantis $D$. Formetur simili modo ut in art.~186 progressio formarum reductarum utrimque infinita, $\mathinner{\ldotp\ldotp\ldotp\ldotp} {}''f,{}'f,f,f',f''\mathinner{\ldotp\ldotp\ldotp\ldotp}$, et quidem sit
\SourceLeaf{177}
\[
\begin{aligned}
f'&=(-a',b',a''),& f''&=(a'',b'',-a''')\quad\text{etc.}\\
{}'f&=(-{}'a,{}'b,a),& {}''f&=({}''a,{}''b,-{}'a)\quad\text{etc.}
\end{aligned}
\]
Ponatur
\[
\begin{array}{llll}
\dfrac{b+b'}{-a'}=h',&\dfrac{b'+b''}{a''}=h'',&\dfrac{b''+b'''}{-a'''}=h'''&\text{etc.}\\[6pt]
\dfrac{{}'b+b}{a}=h,&\dfrac{{}''b+{}'b}{-{}'a}={}'h,&\dfrac{{}'''b+{}''b}{{}''a}={}''h&\text{etc.}
\end{array}
\]
Tum patet, si (ut in art.~177) numeri $\alpha',\alpha'',\alpha'''$ etc. $\beta',\beta'',\beta'''$ etc. etc. formentur secundum algorithmum sequentem
{\small
\setlength{\arraycolsep}{3pt}
\[
\begin{array}{llll}
\alpha'=0&\beta'=-1&\gamma'=1&\delta'=h'\\
\alpha''=\beta'&\beta''=h''\beta'&\gamma''=\delta&\delta''=h''\delta'-1\\
\alpha'''=\beta''&\beta'''=h'''\beta''-\beta'&\gamma'''=\delta''&\delta'''=h'''\delta''-\delta'\\
\alpha''''=\beta'''&\beta''''=h''''\beta'''-\beta''&\gamma''''=\delta'''&\delta''''=h''''\delta'''-\delta''
\\
\multicolumn{4}{c}{\text{etc.}}
\end{array}
\]
}
$f$ transformatum iri
\[
\begin{array}{@{}rlll@{}}
\text{in}&f'&\text{per substitutionem}&\alpha',\ \beta',\ \gamma',\ \delta'\\
&f''&\mathinner{\ldotp\ldotp\ldotp\ldotp\ldotp\ldotp\ldotp\ldotp}&\alpha'',\ \beta'',\ \gamma'',\ \delta''\\
&f'''&\mathinner{\ldotp\ldotp\ldotp\ldotp\ldotp\ldotp\ldotp\ldotp}&\alpha''',\ \beta''',\ \gamma''',\ \delta'''
\end{array}
\]
\centerline{etc.}
omnesque has transformationes fore proprias.

Quum ${}'f$ transeat in $f$ per substitutionem propriam $0,-1,1,h$ (art.~158): $f$ transibit in ${}'f$ per subst. prop. $h,1,-1,0$. Ex simili ratione ${}'f$ transibit in ${}''f$ per subst propr. ${}'h,1,-1,0$; ${}''f$ in ${}'''f$ per subst. pr. ${}''h,1,-1,0$ etc. Hinc per art.~159 eodem modo ut art.~177 colligitur, si numeri ${}'\alpha,{}''\alpha,{}'''\alpha$ etc. ${}'\beta,{}''\beta,{}'''\beta$ etc. etc. formentur secundum algorithmum sequentem
{\small
\setlength{\arraycolsep}{3pt}
\[
\begin{array}{llll}
{}'\alpha=h&{}'\beta=1&{}'\gamma=-1&{}'\delta=0\\
{}''\alpha={}'h\,{}'\alpha-1&{}''\beta={}'\alpha&{}''\gamma={}'h\,{}'\gamma&{}''\delta={}'\gamma\\
{}'''\alpha={}''h\,{}''\alpha-{}'\alpha&{}'''\beta={}''\alpha&{}'''\gamma={}''h\,{}''\gamma-{}'\gamma&{}'''\delta={}''\gamma\\
{}''''\alpha={}'''h\,{}'''\alpha-{}''\alpha&{}''''\beta={}'''\alpha&{}''''\gamma={}'''h\,{}'''\gamma-{}''\gamma&{}''''\delta={}'''\gamma
\\
\multicolumn{4}{c}{\text{etc.}}
\end{array}
\]
}
$f$ transformatum iri
\SourceLeaf{178}
\[
\begin{array}{@{}rlll@{}}
\text{in}&{}'f&\text{per substitutionem}&{}'\alpha,\ {}'\beta,\ {}'\gamma,\ {}'\delta\\
&{}''f&\mathinner{\ldotp\ldotp\ldotp\ldotp\ldotp\ldotp\ldotp\ldotp}&{}''\alpha,\ {}''\beta,\ {}''\gamma,\ {}''\delta\\
&{}'''f&\mathinner{\ldotp\ldotp\ldotp\ldotp\ldotp\ldotp\ldotp\ldotp}&{}'''\alpha,\ {}'''\beta,\ {}'''\gamma,\ {}'''\delta
\end{array}
\]
\centerline{etc.}
omnesque has transformationes fore proprias.

Si ponitur $\alpha=1$, $\beta=0$, $\gamma=0$, $\delta=1$: hi numeri eandem relationem habebunt ad formam $f$, quam habent $\alpha',\beta',\gamma',\delta'$ ad $f'$; $\alpha'',\beta'',\gamma'',\delta''$ ad $f''$ etc.; ${}'\alpha,{}'\beta,{}'\gamma,{}'\delta$ ad ${}'f$ etc. Scilicet per substitutionem $\alpha,\beta,\gamma,\delta$ forma $f$ transibit in $f$. Tum vero progressiones infinitae $\alpha',\alpha'',\alpha'''$ etc., ${}'\alpha,{}''\alpha,{}'''\alpha$ etc., per intercalationem termini $\alpha$, concinne iungentur ita ut unam continuam utrimque infinitam constituere concipi possint secundum eandem legem ubique progredientem $\mathinner{\ldotp\ldotp\ldotp} {}'''\alpha,{}''\alpha,{}'\alpha,\alpha,\alpha',\alpha'',\alpha'''\mathinner{\ldotp\ldotp\ldotp}$ Lex progressionis haec est:
\[
\begin{gathered}
{}'''\alpha+{}'\alpha={}''h\,{}''\alpha,\quad {}''\alpha+\alpha={}'h\,{}'\alpha,\quad {}'\alpha+\alpha'=h\alpha,\\
\alpha+\alpha''=h'\alpha',\quad \alpha'+\alpha'''=h''\alpha''\text{ etc.}
\end{gathered}
\]
sive generaliter (si indicem negativum a dextra scriptum idem designare supponimus, ac positivum a laeva)
\[
\alpha^{m-1}+\alpha^{m+1}=h^m{\alpha^m}
\]
Simili modo progressio $\mathinner{\ldotp\ldotp\ldotp} {}''\beta,{}'\beta,\beta,\beta',\beta''\mathinner{\ldotp\ldotp\ldotp}$ continua erit, cuius lex
\[
\beta^{m-1}+\beta^{m+1}=h^{m+1}\beta^m
\]
et proprie cum praecedente identica, omnibus terminis uno loco promotis, ${}''\beta={}'\alpha$, ${}'\beta=\alpha$, $\beta=\alpha'$ etc. Lex progressionis continuae $\mathinner{\ldotp\ldotp\ldotp\ldotp} {}''\gamma,\allowbreak {}'\gamma,\allowbreak \gamma,\allowbreak \gamma',\allowbreak \gamma''\mathinner{\ldotp\ldotp\ldotp\ldotp\ldotp}$ erit haec
\[
\gamma^{m-1}+\gamma^{m+1}=h^m{\gamma^m}
\]
et lex huius $\mathinner{\ldotp\ldotp\ldotp} {}''\delta,{}'\delta,\delta,\delta',\delta''\mathinner{\ldotp\ldotp\ldotp}$ erit
\[
\delta^{m-1}+\delta^{m+1}=h^{m+1}\delta^m
\]
insuperque generaliter $\delta^m=\gamma^{m+1}$.

\textit{Ex.} Sit forma proposita $f$ haec $(3,8,-5)$, quae transformabitur
\SourceLeaf{179}
\[
\setlength{\arraycolsep}{4pt}
\begin{array}{@{}lll@{}}
\multicolumn{2}{l}{\text{in formam}}&\text{per substitutionem}\\
{}'''''''f&(-10,7,3)&-805,\ -152,\ +143,\ +27\\
{}''''''f&(3,8,-5)&-152,\ +45,\ +27,\ -8\\
{}'''''f&(-5,7,6)&+45,\ +17,\ -8,\ -3\\
{}''''f&(6,5,-9)&+17,\ -11,\ -3,\ +2\\
{}'''f&(-9,4,7)&-11,\ -6,\ +2,\ +1\\
{}''f&(7,3,-10)&-6,\ +5,\ +1,\ -1\\
{}'f&(-10,7,3)&+5,\ +1,\ -1,\ 0\\
f&(3,8,-5)&+1,\ 0,\ 0,\ +1\\
f'&(-5,7,6)&0,\ -1,\ +1,\ -3\\
f''&(6,5,-9)&-1,\ -2,\ -3,\ -7\\
f'''&(-9,4,7)&-2,\ +3,\ -7,\ +10\\
f''''&(7,3,-10)&+3,\ +5,\ +10,\ +17\\
f'''''&(-10,7,3)&+5,\ -8,\ +17,\ -27\\
f''''''&(3,8,-5)&-8,\ -45,\ -27,\ -152\\
f'''''''&(-5,7,6)&-45,\ +143,\ -152,\ +483
\end{array}
\]
\centerline{etc.}

\Needspace{8\baselineskip}
\Article{189}

Circa hunc algorithmum sequentia sunt annotanda:

1) Omnes $a,a',a''$ etc., ${}'a,{}''a$ etc. eadem signa habebunt; omnes $b,b',b''$ etc. ${}'b,{}''b$ etc. erunt positivi; in progressione $\mathinner{\ldotp\ldotp\ldotp} {}''h,{}'h,h,h',h''\mathinner{\ldotp\ldotp\ldotp\ldotp}$ signa alternabunt, scilicet si omnes $a,a'$ etc. sunt positivi, $h^m$ vel ${}^m h$ erit positivus quando $m$ est par, negativus quando $m$ impar; si vero $a,a'$ etc. sunt negativi, $h^m$ vel ${}^m h$ pro $m$ pari erit negativus, pro impari positivus.

2) Si $a$ est positivus adeoque $h'$ negativus, $h''$ positivus etc., erit $\alpha''=-1$ neg., $\alpha'''=h''\alpha''$ neg. et $>\alpha''$ (vel $=\alpha''$ si $h''=1$); $\alpha''''=h'''\alpha'''-\alpha''$ pos. et $>\alpha'''$ (quia $h'''\alpha'''$ pos., $\alpha''$ neg.); $\alpha'''''=h''''\alpha''''-\alpha'''$ pos. et $>\alpha''''$ (quia $h''''\alpha''''$ pos.) etc. Hinc facile concluditur, progressionem $\alpha',\alpha'',\alpha'''$ etc. in infinitum crescere duoque signa positiva semper duo negativa excipere ita ut $\alpha^m$ habeat signum $+,+,-,-$, prout $m\equiv0,1,2,3\ (\mathrm{mod.}\ 4)$. Si $a$ est negativus, per simile ratiocinium invenitur $\alpha''$ neg., $\alpha'''$ pos. et vel $>$ vel $=\alpha''$; $\alpha''''$ pos. $>\alpha'''$; $\alpha'''''$ neg. $>\alpha''''$ etc.,
\SourceLeaf{180}
ita ut progressio $\alpha',\alpha'',\alpha'''$ etc. continuo crescat, signumque termini $\alpha^m$ sit $+,-,-,+$, prout $m\equiv0,1,2,3\ (\mathrm{mod.}\ 4)$.

3) Hoc modo invenitur, omnes quatuor progressiones infinitas $\alpha',\alpha'',\alpha'''$ etc. $\gamma,\gamma',\gamma''$ etc.; $\alpha',\alpha,{}'\alpha,{}''\alpha$ etc.; $\gamma,{}'\gamma,{}''\gamma$ etc. continuo crescere, adeoque etiam sequentes cum illis identicas: $\beta,\beta',\beta''$ etc.; ${}'\delta,\delta,\delta',\delta''$ etc.; $\beta,{}'\beta,{}''\beta$ etc.; ${}'\delta,{}''\delta$ etc.; et, prout $m\equiv0,1,2,3\ (\mathrm{mod.}\ 4)$, signum
\[
\begin{array}{ll}
\text{ipsius }\alpha^m,\quad +\ \pm\ -\ \mp;&\text{ipsius }\beta^m,\quad \pm\ -\ \mp\ +\\[6pt]
\text{ipsius }\gamma^m,\quad \pm\ +\ \mp\ -;&\text{ipsius }\delta^m,\quad +\ \mp\ -\ \pm\\[6pt]
\text{ipsius }{}^m\alpha,\quad +\ \pm\ -\ \mp;&\text{ipsius }{}^m\beta,\quad \mp\ +\ \pm\ -\\[6pt]
\text{ipsius }{}^m\gamma,\quad \mp\ -\ \pm\ +;&\text{ipsius }{}^m\delta,\quad +\ \mp\ -\ \pm
\end{array}
\]
valentibus superioribus quando $a$ est positivus, inferioribus quando $a$ negativus. Teneatur imprimis haec proprietas: Designante $m$ indicem quemcunque positivum, $\alpha^m$ et $\gamma^m$ habebunt eadem signa quando $a$ {positivus. opposita} quando $a$ negativus, similiterque $\beta^m$ et $\delta^m$; contra ${}^m\alpha$ et ${}^m\gamma$, vel ${}^m\beta$ et ${}^m\delta$ habebunt eadem signa quando $a$ negativus, opposita quando $a$ positivus.

4) In signis art. 27 magnitudo ipsorum $\alpha^m$ etc. concinne ita exhiberi potest, ponendo
\[
\begin{gathered}
\mp h'=k',\quad \pm h''=k'',\quad \mp h'''=k'''\text{ etc.}\\
\pm h=k,\quad \mp{}'h={}'k,\quad \pm{}''h={}''k\text{ etc.}
\end{gathered}
\]
ita ut omnes $k',k''$ etc. $k,{}'k$ etc. sint numeri positivi:
{\small
\setlength{\arraycolsep}{3pt}
\[
\begin{array}{ll}
\alpha^m=\pm[k'',k''',k''''\mathinner{\ldotp\ldotp\ldotp\ldotp} k^{m-1}];&\beta^m=\pm[k'',k''',k''''\mathinner{\ldotp\ldotp\ldotp} {k^m]}\\[6pt]
\gamma^m=\pm[k',k'',k'''\mathinner{\ldotp\ldotp\ldotp\ldotp\ldotp} k^{m-1}];&\delta^m=\pm[k',k'',k'''\mathinner{\ldotp\ldotp\ldotp\ldotp} k^m]\\[6pt]
{}^m\alpha=\pm[k,{}'k,{}''k\mathinner{\ldotp\ldotp\ldotp\ldotp\ldotp\ldotp}{}^{m-1}k];&{}^m\beta=\pm[k,{}'k,{}''k\mathinner{\ldotp\ldotp\ldotp\ldotp}{}^{m-2}k]\\[6pt]
{}^m\gamma=\pm[{}'k,{}''k\mathinner{\ldotp\ldotp\ldotp\ldotp\ldotp\ldotp}{}^{m-1}k];&{}^m\delta=\pm[{}'k,{}''k\mathinner{\ldotp\ldotp\ldotp\ldotp}{}^{m-2}k]
\end{array}
\]
}
\textit{signa} vero ad praecepta modo tradita determinari debent. Secundum has formulas, quarum demonstrationem propter facilitatem omittimus, calculus semper expeditissime absolvi poterit.
\SourceLeaf{181}

\Needspace{8\baselineskip}
\Article{190}

\textsc{Lemma.} \textit{Designantibus $m,\mu,m',n,\nu,n'$ numeros integros quoscunque, ita tamen ut trium posteriorum nullus sit $=0$: dico, si $\frac{\mu}{\nu}$ iaceat inter limites $\frac{m}{n}$ et $\frac{m'}{n'}$ exclusive, atque sit $mn'-nm'={\mp1}$, denominatorem $\nu$ fore maiorem quam $n$ et $n'$.}

\textit{Dem.} Manifesto $\mu nn'$ iacebit inter $\nu mn'$ et $\nu nm'$, adeoque ab utroque limite minus differet quam limes alter ab altero, \textit{i. e.} erit $\nu mn'-\nu nm'>\mu nn'-\nu mn'$ et $>\mu nn'-\nu nm'$, sive $\nu>n'(\mu n-\nu m)$ et $>n(\mu n'-\nu m')$. Hinc sequitur, quoniam $\mu n-\nu m$ certe non $=0$ (alioquin enim foret $\frac{\mu}{\nu}=\frac{m}{n}$ contra hyp.), neque $\mu n'-\nu m'=0$ (ex simili ratione), sed uterque ad minimum $=1$, fore $\nu>n'$ et $>n$. \textit{Q. E. D.}

Perspicuum itaque est, $\nu$ non posse esse $=1$, \textit{i. e.} si fuerit $mn'-nm'=\pm1$, inter fractiones $\frac{m}{n}$, $\frac{m'}{n'}$ nullum numerum integrum iacere posse. Quare etiam cifra inter ipsas iacere nequit, \textit{i. e.} fractiones istae signa opposita habere nequeunt.

\Needspace{8\baselineskip}
\Article{191}

\textsc{Theorema.} \textit{Si forma reducta $(a,b,-a')$ determinantis $D$ per substitutionem $\alpha,\beta,\gamma,\delta$ transit in reductam $(A,B,-A')$ eiusdem determinantis: iacebit,} primo, \textit{$\frac{\pm\surd D-b}{a}$ inter $\frac{\alpha}{\gamma}$ et $\frac{\beta}{\delta}$ (siquidem neque $\gamma$ neque $\delta=0$, i. e. si uterque limes est finitus), accepto signo superiori, quando neuter horum limitum habet signum signo ipsius $a$ oppositum (sive clarius, quando aut uterque idem habet, aut alter idem, alter est $=0$), inferiori quando neuter habet idem ut $a$;} secundo \textit{$\frac{\pm\surd D+b}{a'}$ inter $\frac{\gamma}{\alpha}$ et $\frac{\delta}{\beta}$ (siquidem neque $\alpha$ neque $\beta=0$), signo superiori accepto quando limes neuter signum signo ipsius $a'$ (vel $a$) oppositum habet, inferiori quando neuter habet idem ut $a'$}\HistoricalNote{*)}{Manifestum est, alios casus locum habere non posse, quum ex art. praec. propter $\alpha\delta-\beta\gamma=\pm1$, limites bini neque signa opposita {habere;} neque simul $=0$ esse possint.}.

\textit{Dem.} Habentur aequationes
\[
\begin{array}{rcll}
a\alpha\alpha+2b\alpha\gamma-a'\gamma\gamma&=&A&\qquad\mathinner{\ldotp\ldotp\ldotp\ldotp\ldotp\ldotp\ldotp\ldotp}\qquad[1]\\[4pt]
a\beta\beta+2b\beta\delta-a'\delta\delta&=&-A'&\qquad\mathinner{\ldotp\ldotp\ldotp\ldotp\ldotp\ldotp\ldotp}\qquad[2]
\end{array}
\]
Unde deducitur
\[
\begin{array}{rcll}
\dfrac{\alpha}{\gamma}&=&\dfrac{\pm\surd\bigl(D+\frac{aA}{\gamma\gamma}\bigr)-b}{{a}}&\qquad\mathinner{\ldotp\ldotp\ldotp\ldotp\ldotp\ldotp\ldotp\ldotp}\qquad[3]\\[12pt]
\dfrac{\beta}{\delta}&=&\dfrac{\pm\surd\bigl(D-\frac{aA'}{\delta\delta}\bigr)-b}{a}&\qquad\mathinner{\ldotp\ldotp\ldotp\ldotp\ldotp\ldotp\ldotp\ldotp}\qquad[4]
\end{array}
\]
